\begin{frame}
  \frametitle{Outline}
% One \nocite per line -- a multi-key \nocite spanning lines breaks bibtex.
\nocite{Haw71spe}
\nocite{HO74clu}
\nocite{DVJ08int}
\nocite{DZ13exa}
A Hawkes process is a counting process whose events raise the rate of future events.
It is the simplest model of order flow in which a burst is not a coincidence.

We set up the general apparatus --- intensity, compensator, time change ---
fix the exponential kernel that makes it computable,
and read off what the branching ratio does and does not tell us.
\end{frame}

\subsection{Counting processes}

\begin{frame}
  \frametitle{Intensity and compensator}
A multivariate counting process $\multiCountingProc$ has $\filtrationF$-compensator
$\compensator$ when $\multiCountingProc - \compensator$ is a local martingale
and $\compensator$ is predictable, right-continuous and of finite variation.

\pause

When $\compensator(t) = \intzerot \intensity[ ](s)ds$
we call $\intensity[ ]$ the \emph{intensity}, and
\begin{equation*}
 \Expectation\left[ \countingProc(t)-\countingProc(s) \vert \filtrationF_s \right]
 = \Expectation\left[ \int_{s}^{t} \intensity(u) du \vert \filtrationF_s \right] .
\end{equation*}

\pause

Predictability is the whole content: $\intensity(t)$ is a function of the \emph{strict} past.
Every intensity we write is left-continuous for that reason.
\end{frame}

\begin{frame}
  \frametitle{The time change}
\begin{thm}[Meyer, 1971]
 If no two components jump together, the compensator is continuous, and it diverges,
 then rescaling each component's clock by its own compensator produces
 \emph{mutually independent} unit-rate Poisson processes.
\end{thm}

\pause

This is the goodness-of-fit test: rescaled inter-arrivals should be $\text{Exp}(1)$.

\pause

But comparing each margin against $\text{Exp}(1)$ tests only the margins.
The theorem asserts two further independences --- serial, and across components ---
and it is the no-common-jumps hypothesis that buys the second.
\end{frame}

\subsection{Hawkes processes}

\begin{frame}
  \frametitle{The definition}
\begin{equation*}
 \intensity(t) = \baseIntensity_e
 + \sum_{\eone=1}^{\numEventTypes}\intzerot \hawkesKernel\subscriptee(t-s) \, d\countingProc[\eone](s)
\end{equation*}
with $\baseIntensity \geq 0$ and $\hawkesKernel \geq 0$.

\pause

$\branchingMatrix\subscriptee := \Vert \hawkesKernel\subscriptee \Vert_{\Lone}$,
and $\branchingRatio$ its spectral radius: the \emph{branching ratio}.
A stationary version with finite mean intensity exists iff $\branchingRatio<1$.

\pause

\alert{The second index excites.}
$\branchingMatrix\subscriptee$ is the influence of $\eone$ on $e$,
so $\intensity[ ] = \baseIntensity + \excitation\decayedCounts$ is a plain matrix--vector
product and the \emph{column} sums are the readable quantity.
The opposite convention is common, and
$\branchingRatio(\branchingMatrix) = \branchingRatio(\branchingMatrix^{\transpose})$:
a transposed kernel is a bug that still runs.
\end{frame}

\begin{frame}
  \frametitle{Why order flow carries information}
\begin{prop}[Hawkes and Oakes, 1974]
 With $\hawkesKernel \geq 0$ and $\branchingRatio<1$:
 immigrants arrive Poisson at rate $\baseIntensity_e$;
 each event spawns offspring of type $e$ at intensity $\hawkesKernel\subscriptee(s-t)$.
 The superposition is the stationary Hawkes process.
\end{prop}

\pause

Events partition into \emph{clusters}: one immigrant and its descendants.

\pause

A burst is, with high probability, a cluster in progress ---
and a cluster in progress has descendants still to come.
That is the whole reason a window sum of recent flow predicts anything.

\pause

A cluster is a run with a \emph{common ancestor}.
It is not a run of same-signed events.
\end{frame}

\subsection{The exponential kernel}

\begin{frame}
  \frametitle{One decay, and what it buys}
$\hawkesKernel\subscriptee(t) = \excitation\subscriptee e^{-\decay t}$,
with $\decay$ \emph{scalar}.
With
$\decayedCounts_e(t) := \sum_{j:\, \arrivalTimes_{j} < t} e^{-\decay(t-\arrivalTimes_{j})}$,
\begin{equation*}
 \intensity[ ](t) = \baseIntensity + \excitation \decayedCounts(t),
 \qquad
 d\decayedCounts = -\decay \decayedCounts \, dt + d\multiCountingProc .
\end{equation*}

\pause

$(\multiCountingProc, \decayedCounts)$ is Markov: the past enters through $\numEventTypes$
numbers. Updating costs $O(1)$ per event against $O(n^2)$.

\pause

Choosing the kernel is choosing the size of the state.
A decay per pair needs a $\numEventTypes \times \numEventTypes$ state.
The price of the single timescale is that the model cannot have two.
\end{frame}

\begin{frame}
  \frametitle{Exact simulation}
The decay is common, so the excited part of $\totalIntensity$ decays by one factor:
\begin{equation*}
 \Prob \left( \arrivalTimes[]_{n+1} - \arrivalTimes[]_{n} > s \, \vert \, \filtrationF \right)
 = e^{-\bar{\baseIntensity}s} \cdot
 \exp\left( -\tfrac{D}{\decay}\big(1 - e^{-\decay s}\big) \right) .
\end{equation*}

\pause

A product of two survival functions, so the wait is $\tau_1 \wedge \tau_2$:
one immigration draw, one excitation draw, both closed form.
No thinning, no rejection \citep{DZ13exa}.

\pause

$\tau_2 = +\infty$ is not an edge case:
the excited part carries finite mass $D/\decay$ and expires without firing with probability
$e^{-D/\decay}$.

\pause

The \emph{type} is then drawn proportionally to $\intensity(\arrivalTimes[]_{n+1}-)$.
The wait sees only column sums; the type draw is where $\excitation$ re-enters.
\end{frame}

\subsection{Reading the branching ratio}

\begin{frame}
  \frametitle{What $\branchingRatio$ is, and what it is not}
$\stationaryIntensity = (I-\branchingMatrix)^{-1}\baseIntensity$,
$\totalRate = \mathbf{1}^{\transpose}\stationaryIntensity$,
$\bar{\baseIntensity} = \mathbf{1}^{\transpose}\baseIntensity$.

\pause

Cluster size $C_j = \mathbf{1}^{\transpose}(I-\branchingMatrix)^{-1}e_j$
--- \alert{counting the ancestor} --- and
$\bar{C} = \totalRate/\bar{\baseIntensity}$, so $\bar{\baseIntensity}\bar{C} = \totalRate$:
every event lies in exactly one cluster.

\pause

Endogenous fraction $1 - \bar{\baseIntensity}/\totalRate$.

\pause

\alert{It is not $\branchingRatio$}, and $\bar C$ is not $1/(1-\branchingRatio)$,
except when the column sums of $\branchingMatrix$ are constant
(or $\baseIntensity$ is a right Perron vector).
Here: $0.682$ against $\branchingRatio = 0.60$, and $\bar C = 3.15$ against $2.5$.
A spectral radius does not know about $\baseIntensity$.
\end{frame}

\begin{frame}
  \frametitle{Timescales, and how to switch a study off}
Relaxation $1/(\decay(1-\branchingRatio))$ --- the \emph{slowest mode}, not a half-life.

\pause

Two dimensionless groups:
\begin{itemize}
 \item $\totalRate/\decay$: events per excitation timescale. Is clustering visible at all?
 \item $\totalRate w$: events per observation window. Is a window sum a statistic?
\end{itemize}

\pause

At $\totalRate/\decay \ll 1$ the excitation has expired before the next event arrives.
A study run there measures a process whose self-excitation is switched off,
and correctly finds nothing.
\end{frame}

\subsection{Reading the order flow}

\begin{frame}
  \frametitle{What the laboratory leaves out}
In a traded market some flow is \emph{informed}.
A provider resting on the bid is short an option:
he is filled exactly when someone wants to sell, which is disproportionately when selling
is right.

\pause

His defence is to withdraw first.
So a queue about to be picked off is one whose own side pulls away from it:
\alert{depletion begets depletion}.
That is what makes $\queueImbalance[1]$ informative --- a thin bid is thin \emph{because}
its providers inferred something.

\pause

Our generator does the reverse, deliberately: depletion excites \emph{replenishment},
19 to 1.
The two cannot both be strong in one kernel, and a book that survives is the precondition
for measuring anything.

\pause

So here $\queueImbalance[1]$ carries \alert{mechanical} information only ---
a depleted queue is one market order from clearing.
Read the comparison asymmetrically:
$\OFI$ winning is not evidence it dominates in general;
$\queueImbalance[1]$ winning would be the stronger finding.
\end{frame}
